(3 points + 3 points)
\begin{teilaufgaben}
\item
The function $u(x)$ is defined on the interval $\Omega = [-2,2].$ 

The function $u(x)$ satisfies in $\Omega$ 
\[
u''(x) = u(x)
\]
and  $u(-2) = 7, \ u(2) = 7$ on the boundary of $\Omega$. 

Determine approximate values for $u(-1)$, $u(0)$ and $u(1)$. 

Use centered finite differences with $h = 1$. 

\item
The function $u(x,y)$ is defined on the square $\Omega = [0, 3] \times [0,3]$.

The function $u(x,y)$ satisfies in $\Omega$
\[
\Delta u(x,y) = u(x,y)
\]

and $u(x,y) = 3$ on the boundary of $\Omega$.

Determine approximate values for $u(1,1)$, $u(2,1)$, $u(1,2)$, $u(2,2)$.

Use finite differences with $h = 1.$

\end{teilaufgaben}

\begin{loesung}
\begin{teilaufgaben}
\item
By the method of finite differences we have
\begin{equation*}
\renewcommand{\arraycolsep}{2.4pt}
\begin{array}{lcrclclcl}
\tilde u(-2) &-& 3 &\cdot& \tilde u(-1) &+& \tilde u(0) &=& 0, \\
\tilde u(-1) &-& 3 &\cdot& \tilde u(0)  &+& \tilde u(1) &=& 0, \\
\tilde u(0)  &-& 3 &\cdot& \tilde u(1)  &+& \tilde u(2) &=& 0.
\end{array}
\end{equation*}
%\[
%\tilde u(-2) - 3 \cdot \tilde u(-1) + \tilde u(0) = 0,
%\]
%\[
%\tilde u(-1) - 3 \cdot \tilde u(0) + \tilde u(1) = 0,
%\]
%\[
%\tilde u(0) - 3 \cdot \tilde u(1) + \tilde u(2) = 0.
%%\]
This leads to a system of three linear equations:
\begin{equation}
\begin{pmatrix*}[r]
-3 & 1 & 0 \\
 1 & -3 & 1\\
 0 & 1 & -3
\end{pmatrix*}
\cdot
\begin{pmatrix}
\tilde u(-1)  \\
  \tilde u(0) \\
\tilde u(1)
\end{pmatrix}
=
\begin{pmatrix*}[r]
-7  \\
 0 \\
-7
\end{pmatrix*}.
\tag{\textbf{2P}}
\end{equation}

Its solution is
\begin{equation}
(\tilde u(-1),  \tilde u(0), \tilde u(1)) = (3 , 2, 3).
\tag{\textbf{1P}}
\end{equation}

\item
By the five-point-star operator we have the following system of
four linear equations:
\begin{equation}
\begin{pmatrix*}[r]
-5 &  1 &  1 &  0 \\
 1 & -5 &  0 &  1 \\
 1 &  0 & -5 &  1 \\
 0 &  1 &  1 & -5
\end{pmatrix*}
\cdot
\begin{pmatrix}
\tilde u(1,1) \\
\tilde u(2,1) \\
\tilde u(1,2) \\
\tilde u(2,2)
\end{pmatrix}
+
\begin{pmatrix*}[r] 6 \\  6 \\ 6 \\ 6 \end{pmatrix*}
=
\begin{pmatrix}
\tilde u(1,1) \\
\tilde u(2,1) \\
\tilde u(1,2) \\
\tilde u(2,2)
\end{pmatrix}.
\tag{\textbf{2P}}
\end{equation}
Hence
\begin{equation}
\tilde u(1,1)
=
\tilde u(1,2)
=
\tilde u(2,1)
=
\tilde u(2,2)
=
2.
\tag{\textbf{1P}}
\end{equation}
\end{teilaufgaben}
\end{loesung}
